% uid: ""
% school: "lycee"
% level: "1ere"
% chapter: "Activités numériques"
% title: "Restes"
% difficulty: "Moyen"
% tags: ["division-euclidienne", "restes", "divisibilite", "entiers"]
% provenance: "own"
% batch: "2026-11-22-revision-1ere-A"

% ---STATEMENT---
\begin{enumerate}
  \item Montrer que le carré d'un entier impair a pour reste $1$ dans la division euclidienne par $8$. [Écrire $n=8q+r$ avec $r\in\{1,3,5,7\}$.]
  \item Un entier $n$ a pour reste $7$ dans la division euclidienne par $25$. Quels sont les restes de $n^2$ et de $n^3$ dans la division par $25$ ?
  \item Un entier naturel $n$ a pour reste $3$ dans la division par $7$ et pour reste $2$ dans la division par $5$. Montrer que $n-17$ est un multiple de $35$, en déduire l'écriture de $n$, puis le plus petit entier $n>2026$ qui convient.
\end{enumerate}

% ---SOLUTION---
\begin{prop}[Division euclidienne]
Pour deux entiers $a$ et $b>0$, il existe un unique couple $(q,r)$ tel que \ib{$a=bq+r$ avec $0\leq r<b$}. Si deux entiers \textbf{premiers entre eux} divisent $m$, leur produit divise $m$.
\end{prop}
\begin{enumerate}
  \item Un entier impair a pour reste $r\in\{1,3,5,7\}$ dans la division par $8$ : $n=8q+r$.
  \[n^2=64q^2+16qr+r^2=8\,(8q^2+2qr)+r^2\]
  Il reste à étudier $r^2$ :
  \[\begin{array}{c|c|c|c|c}
  r&1&3&5&7\\\hline
  r^2&1&9=8\times1+1&25=8\times3+1&49=8\times6+1
  \end{array}\]
  Dans chaque cas $r^2=8m+1$, donc $n^2=8\,(8q^2+2qr+m)+1$ : le reste est \ib{$1$}.

  \item On écrit $n=25q+7$.

  \textbf{Pour $n^2$ :}
  \[n^2=625q^2+350q+49=25\,(25q^2+14q)+49\]
  Comme $49=25\times1+24$ :
  \[n^2=25\,(25q^2+14q+1)+24\]
  Le reste est \ib{$24$}.

  \textbf{Pour $n^3$ :} on développe $(25q+7)^3$. Tous les termes contiennent le facteur $25q$, sauf le dernier :
  \[n^3=(25q)^3+3(25q)^2\times7+3\times25q\times49+343=25\times(\cdots)+343\]
  Comme $343=13\times25+18$, le reste est \ib{$18$}.

  \item On écrit $n=7a+3=5b+2$ avec $a,b$ entiers. On remarque que $17=7\times2+3=5\times3+2$.

  On soustrait :
  \[n-17=7a+3-17=7\,(a-2)\qquad n-17=5b+2-17=5\,(b-3)\]
  Donc $n-17$ est multiple de $7$ \textbf{et} de $5$. Comme $5$ et $7$ sont premiers entre eux, $n-17$ est multiple de $35$ :
  \[n-17=35k\implies\ib{$n=35k+17$}\ \ (k\in\mathbb N)\]
  On cherche $n>2026$ :
  \[35k+17>2026\iff35k>2009\iff k>57{,}4\]
  Le plus petit entier est $k=58$ :
  \[n=35\times58+17=2030+17=\ib{$2047$}\]
  Contrôle : $2047=7\times292+3$ et $2047=5\times409+2$.
\end{enumerate}
\textbf{Idée :} remplacer $n$ par $n-17$ transforme deux conditions \og reste 3 \fg{} et \og reste 2 \fg{} en deux conditions \og multiple de \fg{}.

\textbf{Piège :} conclure \og multiple de $35$ \fg{} sans avoir vérifié que $5$ et $7$ sont premiers entre eux.
